\subsection{\normalsize Simuladores cuánticos}
La simulación clásica de circuitos cuánticos es esencial para el desarrollo y
validación de algoritmos cuánticos. Sin embargo, el principal desafío radica
en la representación del estado cuántico, cuyo tamaño crece exponencialmente con
el número de qubits.

Uno de los primeros enfoques para abordar este problema fue la representación
mediante \textit{Matrix Product States} (MPS), propuesta por \textcite{Vidal2003}.
Este método permite representar ciertos estados cuánticos de forma eficiente
cuando el entrelazamiento es limitado, reduciendo los requisitos de almacenamiento
y cómputo. Su principal ventaja es la escalabilidad en circuitos con baja
profundidad, permitiendo simular cientos de qubits. No obstante, su limitación
radica en que para sistemas altamente entrelazados el costo computacional crece
exponencialmente, restringiendo su aplicabilidad a ciertos algoritmos.

A medida que se buscaban métodos más generales, surgió la simulación tipo
\textit{Schrödinger}, que almacena el vector de estado completo y lo actualiza
en cada operación cuántica. Este enfoque, empleado en simuladores como QuEST y
Qiskit Aer \parencite{Jones2019, qiskit_aer_doc}, garantiza alta precisión y permite simular cualquier
circuito cuántico. Su desventaja principal es la extrema demanda de memoria, la
cual crece exponencialmente con el número de qubits: una simulación de 45 qubits
ya ha requerido 0.5 petabytes de memoria \parencite{Haner2017}, y al extender esta
representación de doble precisión a 50 qubits el costo pasa al orden de decenas de
petabytes.

Para mitigar el problema de memoria, se desarrollaron métodos basados en
partición del circuito y caminos de Feynman. Por ejemplo, Chen y cols.\
presentaron un algoritmo distribuido capaz
de calcular amplitudes de salida y simular circuitos universales de hasta
$12\times12$ qubits con profundidad moderada sin materializar el estado completo
\parencite{Chen2018}. Sin embargo, aunque reducen el
almacenamiento requerido, estos métodos incrementan la complejidad computacional,
ya que el número de caminos crece exponencialmente con la profundidad del
circuito.
